% Generated from docs/book.pdf. Review equations and figures against the source PDF.
\chapter{Bayesian Decision Theory}
\label{chap:03-bayesian-decision-theory}

We discuss probability theory as the framework for making decisions under uncertainty. In classification, Bayes’ rule is used to calculate the probabilities of the classes. We generalize to discuss how we can make rational decisions among multiple actions to minimize expected risk. We also discuss learning association rules from data.

\section{Introduction}

Programming computers to make inference from data is a cross between statistics and computer science, where statisticians provide the mathematical framework of making inference from data and computer scientists work on the efficient implementation of the inference methods. Data comes from a process that is not completely known. This lack of knowledge is indicated by modeling the process as a random process. Maybe the process is actually deterministic, but because we do not have access to complete knowledge about it, we model it as random and use probability theory to analyze it. At this point, it may be a good idea to jump to the appendix and review basic probability theory before continuing with this chapter. Tossing a coin is a random process because we cannot predict at any toss whether the outcome will be heads or tails---that is why we toss coins, or buy lottery tickets, or get insurance. We can only talk about the probability that the outcome of the next toss will be heads or tails. It may be argued that if we have access to extra knowledge such as the exact composition of the coin, its initial position, the force and its direction that is applied to the coin when tossing it, where and how it is caught, and so forth, the exact outcome of the toss can be predicted.

The extra pieces of knowledge that we do not have access to are named unobservable the unobservable variables. In the coin tossing example, the only observvariables able variable is the outcome of the toss. Denoting the unobservables by observable variable z and the observable as x, in reality we have

x = f(z)

where f($\cdot$) is the deterministic function that defines the outcome from the unobservable pieces of knowledge. Because we cannot model the process this way, we define the outcome X as a random variable drawn from a probability distribution P(X = x) that specifies the process. The outcome of tossing a coin is heads or tails, and we define a random variable that takes one of two values. Let us say X = 1 denotes that the outcome of a toss is heads and X = 0 denotes tails. Such X are Bernoullidistributed where the parameter of the distribution po is the probability that the outcome is heads:

P(X = 1) = po and P(X = 0) = 1 $-$P(X = 1) = 1 $-$po

Assume that we are asked to predict the outcome of the next toss. If we know po, our prediction will be heads if po > 0.5 and tails otherwise. This is because if we choose the more probable case, the probability of error, which is 1 minus the probability of our choice, will be minimum. If this is a fair coin with po = 0.5, we have no better means of prediction than choosing heads all the time or tossing a fair coin ourselves! If we do not know P(X) and want to estimate this from a given sample, sample then we are in the realm of statistics. We have a sample, X, containing examples drawn from the probability distribution of the observables xt, denoted as p(x). The aim is to build an approximator to it, \textasciicircum{}p(x), using the sample X. In the coin tossing example, the sample contains the outcomes of the past N tosses. Then using X, we can estimate po, which is the parameter that uniquely specifies the distribution. Our estimate of po is \#\{tosses with outcome heads\} \textasciicircum{}po = \#\{tosses\} Numerically using the random variables, xt is 1 if the outcome of toss t is heads and 0 otherwise. Given the sample \{heads, heads, heads, tails, heads, tails, tails, heads, heads\}, we have X = \{1, 1, 1, 0, 1, 0, 0, 1, 1\} and the estimate is N t=1 xt 6 \textasciicircum{}po = = N 9

\section{Classification}

We discussed credit scoring in section 1.2.2, where we saw that in a bank, according to their past transactions, some customers are low-risk in that they paid back their loans and the bank profited from them and other customers are high-risk in that they defaulted. Analyzing this data, we would like to learn the class “high-risk customer” so that in the future, when there is a new application for a loan, we can check whether that person obeys the class description or not and thus accept or reject the application. Using our knowledge of the application, let us say that we decide that there are two pieces of information that are observable. We observe them because we have reason to believe that they give us an idea about the credibility of a customer. Let us say, for example, we observe customer’s yearly income and savings, which we represent by two random variables X1 and X2. It may again be claimed that if we had access to other pieces of knowledge such as the state of economy in full detail and full knowledge about the customer, his or her intention, moral codes, and so forth, whether someone is a low-risk or high-risk customer could have been deterministically calculated. But these are nonobservables and with what we can observe, the credibility of a customer is denoted by a Bernoulli random variable C conditioned on the observables X = [X1, X2]T where C = 1 indicates a high-risk customer and C = 0 indicates a low-risk customer. Thus if we know P(C|X1, X2), when a new application arrives with X1 = x1 and X2 = x2, we can

C = 1 if P(C = 1|x1, x2) > 0.5 choose C = 0 otherwise

or equivalently

C = 1 if P(C = 1|x1, x2) > P(C = 0|x1, x2) (3.1) choose C = 0 otherwise

The probability of error is 1 $-$max(P(C = 1|x1, x2), P(C = 0|x1, x2)). This example is similar to the coin tossing example except that here, the Bernoulli random variable C is conditioned on two other observable variables. Let us denote by x the vector of observed variables, x = [x1, x2]T. Bayes’ rule The problem then is to be able to calculate P(C|x). Using Bayes’ rule, it can be written as P(C)p(x|C) (3.2) P(C|x) = p(x)

prior probability P(C = 1) is called the prior probability that C takes the value 1, which in our example corresponds to the probability that a customer is highrisk, regardless of the x value. It is called the prior probability because it is the knowledge we have as to the value of C before looking at the observables x, satisfying

P(C = 0) + P(C = 1) = 1

class likelihood p(x|C) is called the class likelihood and is the conditional probability that an event belonging to C has the associated observation value x. In our case, p(x1, x2|C = 1) is the probability that a high-risk customer has his or her X1 = x1 and X2 = x2. It is what the data tells us regarding the class. evidence p(x), the evidence, is the marginal probability that an observation x is seen, regardless of whether it is a positive or negative example.

(3.3) p(x) = p(x, C) = p(x|C = 1)P(C = 1) + p(x|C = 0)P(C = 0)

C

Combining the prior and what the data tells us using Bayes’ rule, we posterior calculate the posterior probability of the concept, P(C|x), after having probability seen the observation, x.

prior $\times$ likelihood posterior = evidence Because of normalization by the evidence, the posteriors sum up to 1:

P(C = 0|x) + P(C = 1|x) = 1

Once we have the posteriors, we decide by using equation 3.1. For now, we assume that we know the prior and likelihoods; in later chapters, we discuss how to estimate P(C) and p(x|C) from a given training sample. In the general case, we have K mutually exclusive and exhaustive classes; Ci, i = 1, . . . , K; for example, in optical digit recognition, the input is a bitmap image and there are ten classes. We have the prior probabilities satisfying

K (3.4) P(Ci) $\geq$0 and P(Ci) = 1 i=1

p(x|Ci) is the probability of seeing x as the input when it is known to belong to class Ci. The posterior probability of class Ci can be calculated as p(x|Ci)P(Ci) p(x|Ci)P(Ci) (3.5) P(Ci|x) = = K p(x) k=1 p(x|Ck)P(Ck)

Bayes’ classifier and for minimum error, the Bayes’ classifier chooses the class with the highest posterior probability; that is, we

(3.6) choose Ci if P(Ci|x) = max P(Ck|x) k

\section{Losses and Risks}

It may be the case that decisions are not equally good or costly. A financial institution when making a decision for a loan applicant should take into account the potential gain and loss as well. An accepted low-risk applicant increases profit, while a rejected high-risk applicant decreases loss. The loss for a high-risk applicant erroneously accepted may be different from the potential gain for an erroneously rejected low-risk applicant. The situation is much more critical and far from symmetry in other domains like medical diagnosis or earthquake prediction. Let us define action $\alpha$i as the decision to assign the input to class Ci loss function and $\lambda$ik as the loss incurred for taking action $\alpha$i when the input actually expected risk belongs to Ck. Then the expected risk for taking action $\alpha$i is

K (3.7) R($\alpha$i|x) = $\lambda$ikP(Ck|x) k=1

and we choose the action with minimum risk:

(3.8) choose $\alpha$i if R($\alpha$i|x) = min R($\alpha$k|x) k

Let us define K actions $\alpha$i, i = 1, . . . , K, where $\alpha$i is the action of assign- 0/1 loss ing x to Ci. In the special case of the 0/1 loss case where

0 if i = k (3.9) $\lambda$ik = 1 if i /= k

all correct decisions have no loss and all errors are equally costly. The risk of taking action $\alpha$i is

K R($\alpha$i|x) = $\lambda$ikP(Ck|x) k=1

= P(Ck|x) k/=i = 1 $-$P(Ci|x)

because k P(Ck|x) = 1. Thus to minimize risk, we choose the most probable case. In later chapters, for simplicity, we will always assume this case and choose the class with the highest posterior, but note that this is indeed a special case and rarely do applications have a symmetric, 0/1 loss. In the general case, it is a simple postprocessing to go from posteriors to risks and to take the action to minimize the risk. In some applications, wrong decisions---namely, misclassifications--- may have very high cost, and it is generally required that a more complex--- for example, manual---decision is made if the automatic system has low certainty of its decision. For example, if we are using an optical digit recognizer to read postal codes on envelopes, wrongly recognizing the code causes the envelope to be sent to a wrong destination. reject In such a case, we define an additional action of reject or doubt, $\alpha$K+1, with $\alpha$i, i = 1, . . . , K, being the usual actions of deciding on classes Ci, i = 1, . . . , K (Duda, Hart, and Stork 2001). A possible loss function is ( |( 0 if i = k (3.10) $\lambda$ik = $\lambda$ if i = K + 1 |( 1 otherwise

where 0 < $\lambda$ < 1 is the loss incurred for choosing the (K + 1)st action of reject. Then the risk of reject is

K (3.11) R($\alpha$K+1|x) = $\lambda$P(Ck|x) = $\lambda$ k=1

and the risk of choosing class Ci is

(3.12) R($\alpha$i|x) = P(Ck|x) = 1 $-$P(Ci|x) k/=i

The optimal decision rule is to

choose Ci if R($\alpha$i|x) < R($\alpha$k|x) for all k /= i and

R($\alpha$i|x) < R($\alpha$K+1|x)

(3.13) reject if R($\alpha$K+1|x) < R($\alpha$i|x), i = 1, . . . , K

Given the loss function of equation 3.10, this simplifies to

choose Ci if P(Ci|x) > P(Ck|x) for all k /= i and

P(Ci|x) > 1 $-$$\lambda$

(3.14) reject otherwise

This whole approach is meaningful if 0 < $\lambda$ < 1: If $\lambda$ = 0, we always reject; a reject is as good as a correct classification. If $\lambda$ $\geq$1, we never reject; a reject is as costly as, or costlier than, an error.

\section{Discriminant Functions}

discriminant Classification can also be seen as implementing a set of discriminant funcfunctions tions, gi(x), i = 1, . . . , K, such that we

(3.15) choose Ci if gi(x) = max gk(x) k

We can represent the Bayes’ classifier in this way by setting

gi(x) = $-$R($\alpha$i|x)

and the maximum discriminant function corresponds to minimum conditional risk. When we use the 0/1 loss function, we have

gi(x) = P(Ci|x)

or ignoring the common normalizing term, p(x), we can write

gi(x) = p(x|Ci)P(Ci)

decision regions This divides the feature space into K decision regions R1, . . . , RK, where Ri = \{x|gi(x) = maxk gk(x)\}. The regions are separated by decision boundaries, surfaces in feature space where ties occur among the largest discriminant functions (see figure 3.1). When there are two classes, we can define a single discriminant

g(x) = g1(x) $-$g2(x)

and we

C1 if g(x) > 0 choose C2 otherwise

An example is a two-class learning problem where the positive examples can be taken as C1 and the negative examples as C2. When K = 2, dichotomizer the classification system is a dichotomizer and for K $\geq$3, it is a polypolychotomizer chotomizer.

x2 C1

reject

C2 C3

x1

Figure 3.1 Example of decision regions and decision boundaries.

\section{Utility Theory}

In equation 3.7, we defined the expected risk and chose the action that utility theory minimizes expected risk. We now generalize this to utility theory, which is concerned with making rational decisions when we are uncertain about the state. Let us say that given evidence x, the probability of state Sk is utility function calculated as P(Sk|x). We define a utility function, Uik, which measures expected utility how good it is to take action $\alpha$i when the state is Sk. The expected utility is

(3.16) EU($\alpha$i|x) = UikP(Sk|x) k

A rational decision maker chooses the action that maximizes the expected utility

(3.17) Choose $\alpha$i if EU($\alpha$i|x) = max EU($\alpha$j|x) j

In the context of classification, decisions correspond to choosing one of the classes, and maximizing the expected utility is equivalent to minimizing expected risk. Uik are generally measured in monetary terms, and this gives us a way to define the loss matrix $\lambda$ik as well. For example, in

defining a reject option (equation 3.10), if we know how much money we will gain as a result of a correct decision, how much money we will lose on a wrong decision, and how costly it is to defer the decision to a human expert, depending on the particular application we have, we can fill in the correct values Uik in a currency unit, instead of 0, $\lambda$, and 1, and make our decision so as to maximize expected earnings. Note that maximizing expected utility is just one possibility; one may define other types of rational behavior, for example, minimizing the worst possible loss. In the case of reject, we are choosing between the automatic decision made by the computer program and human decision that is costlier but assumed to have a higher probability of being correct. Similarly one can imagine a cascade of multiple automatic decision makers, which as we proceed are costlier but have a higher chance of being correct; we are going to discuss such cascades in chapter 17 where we talk about combining multiple learners.

\section{Association Rules}

association rule An association rule is an implication of the form X $\to$Y where X is the antecedent and Y is the consequent of the rule. One example of associbasket analysis ation rules is in basket analysis where we want to find the dependency between two items X and Y. The typical application is in retail where X and Y are items sold, as we discussed in section 1.2.1. In learning association rules, there are three measures that are frequently calculated:

support Support of the association rule X $\to$Y:

\#\{customers who bought X and Y\} (3.18) Support(X, Y) $\equiv$P(X, Y) = \#\{customers\}

confidence Confidence of the association rule X $\to$Y:

P(X, Y) Confidence(X $\to$Y) $\equiv$P(Y|X) = P(X) \#\{customers who bought X and Y\} (3.19) = \#\{customers who bought X\}

lift Lift, also known as interest of the association rule X $\to$Y: interest

P(X, Y) P(Y|X) (3.20) Lift(X $\to$Y) = = P(X)P(Y) P(Y)

There are other measures as well (Omiecinski 2003), but these three, especially the first two, are the most widely known and used. Confidence is the conditional probability, P(Y|X), which is what we normally calculate. To be able to say that the rule holds with enough confidence, this value should be close to 1 and significantly larger than P(Y), the overall probability of people buying Y. We are also interested in maximizing the support of the rule, because even if there is a dependency with a strong confidence value, if the number of such customers is small, the rule is worthless. Support shows the statistical significance of the rule, whereas confidence shows the strength of the rule. The minimum support and confidence values are set by the company, and all rules with higher support and confidence are searched for in the database. If X and Y are independent, then we expect lift to be close to 1; if the ratio differs---if P(Y|X) and P(Y) are different---we expect there to be a dependency between the two items: If the lift is more than 1, we can say that X makes Y more likely, and if the lift is less than 1, having X makes Y less likely. These formulas can easily be generalized to more than two items. For example, \{X, Y, Z\} is a three-item set, and we may look for a rule, such as X, Z $\to$Y, that is, P(Y|X, Z). We are interested in finding all such rules having high enough support and confidence and because a sales database is generally very large, we want to find them by doing a small number of Apriori algorithm passes over the database. There is an efficient algorithm, called Apriori (Agrawal et al. 1996) that does this, which has two steps: (1) finding frequent itemsets, that is, those which have enough support, and (2) converting them to rules with enough confidence, by splitting the items into two, as items in the antecedent and items in the consequent:

1. To find frequent itemsets quickly (without complete enumeration of all subsets of items), the Apriori algorithm uses the fact that for \{X, Y, Z\} to be frequent (have enough support), all its subsets \{X, Y\}, \{X, Z\}, and \{Y, Z\} should be frequent as well---adding another item can never increase support. That is, we only need to check for three-item sets all of whose two-item subsets are frequent; or, in other words, if a twoitem set is known not to be frequent, all its supersets can be pruned and need not be checked.

We start by finding the frequent one-item sets and at each step, inductively, from frequent k-item sets, we generate candidate k+1-item sets and then do a pass over the data to check if they have enough support. The Apriori algorithm stores the frequent itemsets in a hash table for easy access. Note that the number of candidate itemsets will decrease very rapidly as k increases. If the largest itemset has n items, we need a total of n + 1 passes over the data.

2. Once we find the frequent k-item sets, we need to convert them to rules by splitting the k items into two as antecedent and consequent. Just like we do for generating the itemsets, we start by putting a single consequent and k $-$1 items in the antecedent. Then, for all possible single consequents, we check if the rule has enough confidence and remove it if it does not.

Note that for the same itemset, there may be multiple rules with different subsets as antecedent and consequent. Then, inductively, we check whether we can move another item from the antecedent to the consequent. Rules with more items in the consequent are more specific and more useful. Here, as in itemset generation, we use the fact that to be able to have rules with two items in the consequent with enough confidence, each of the two rules with single consequent by itself should have enough confidence; that is, we go from one consequent rules to two consequent rules and need not check for all possible two-term consequents (exercise 7).

It should be kept in mind that a rule X $\to$Y need not imply causality hidden variables but just an association. In a problem, there may also be hidden variables whose values are never known through evidence. The advantage of using hidden variables is that the dependency structure can be more easily defined. For example, in basket analysis when we want to find the dependencies among items sold, let us say we know that there is a dependency among “baby food,” “diapers,” and “milk” in that a customer buying one of these is very much likely to buy the other two. Instead of representing dependencies among these three, we may designate a hidden node, “baby at home,” as the hidden cause of the consumption of these three items. Graphical models that we will discuss in chapter 16 allow us to represent such hidden variables. When there are hidden nodes, their values are estimated given the values of observed nodes and filled in.

\section{Notes}

Making decisions under uncertainty has a long history, and over time humanity has looked at all sorts of strange places for evidence to remove the uncertainty: stars, crystal balls, and coffee cups. Reasoning from meaningful evidence using probability theory is only a few hundred years old; see Newman 1988 for the history of probability and statistics and some very early articles by Laplace, Bernoulli, and others who have founded the theory. Russell and Norvig (1995) give an excellent discussion of utility theory and the value of information, also discussing the assignment of utilities in monetary terms. Shafer and Pearl 1990 is an early collection of articles on reasoning under uncertainty. Association rules are successfully used in many data mining applications, and we see such rules on many Web sites that recommend books, movies, music, and so on. The algorithm is very simple and its efficient implementation on very large databases is critical (Zhang and Zhang 2002; Li 2006). Later, we will see in chapter 16 about graphical models how to generalize from association rules to concepts that need not be binary and where associations can be of different types, also allowing hidden variables.

\section{Exercises}

likelihood ratio 1. In a two-class problem, the likelihood ratio is p(x|C1) p(x|C2) Write the discriminant function in terms of the likelihood ratio.

log odds 2. In a two-class problem, the log odds is defined as P(C1|x) log P(C2|x) Write the discriminant function in terms of the log odds.

3. In a two-class, two-action problem, if the loss function is $\lambda$11 = $\lambda$22 = 0, $\lambda$12 = 10, and $\lambda$21 = 1, write the optimal decision rule.

4. Propose a three-level cascade where when one level rejects, the next one is used as in equation 3.10. How can we fix the $\lambda$ on different levels?

5. Somebody tosses a fair coin and if the result is heads, you get nothing; otherwise you get \$5. How much would you pay to play this game? What if the win is \$500 instead of \$5?

6. Generalize the confidence and support formulas for basket analysis to calculate k-dependencies, namely, P(Y|X1, . . ., Xk).

7. Show that as we move an item from the consequent to the antecedent, confidence can never increase: confidence(ABC $\to$D) $\geq$confidence(AB $\to$CD).

8. Associated with each item sold in basket analysis, if we also have a number indicating how much the customer enjoyed the product, for example, in a scale of 0 to 10, how can you use this extra information to calculate which item to propose to a customer?

9. Show example transaction data where for the rule X $\to$Y:

(a) Both support and confidence are high.

(b) Support is high and confidence is low.

(c) Support is low and confidence is high.

(d) Both support and confidence are low.

\section{References}

Agrawal, R., H. Mannila, R. Srikant, H. Toivonen, and A. Verkamo. 1996. “Fast Discovery of Association Rules.” In Advances in Knowledge Discovery and Data Mining, ed. U. M. Fayyad, G. Piatetsky-Shapiro, P. Smyth, and R. Uthurusamy, 307--328. Cambridge, MA: MIT Press.

Duda, R. O., P. E. Hart, and D. G. Stork. 2001. Pattern Classification, 2nd ed. New York: Wiley.

Li, J. 2006. “On Optimal Rule Discovery.” IEEE Transactions on Knowledge and Data Discovery 18: 460--471.

Newman, J. R., ed. 1988. The World of Mathematics. Redmond, WA: Tempus.

Omiecinski, E. R. 2003. “Alternative Interest Measures for Mining Associations in Databases.” IEEE Transactions on Knowledge and Data Discovery 15: 57--69.

Russell, S., and P. Norvig. 1995. Artificial Intelligence: A Modern Approach. New York: Prentice Hall.

Shafer, G., and J. Pearl, eds. 1990. Readings in Uncertain Reasoning. San Mateo, CA: Morgan Kaufmann.

Zhang, C., and S. Zhang. 2002. Association Rule Mining: Models and Algorithms. New York: Springer.
